\section{The Controlling Variable Is the Seed Fraction}
\label{sec:regime}

\subsection{Why absolute budgets mislead}

Our first two attempts to characterise this problem produced the conclusion that out-degree
\emph{does not} fail under overexposure, and both were wrong for the same reason. We sampled
candidate nodes uniformly from the whole graph and correlated out-degree with the true marginal
gain. On a heavy-tailed graph such a pool consists almost entirely of low-degree nodes, so the
correlation is dominated by the global regularity ``high-degree nodes are more valuable'' rather
than by any ability to rank candidates under a given state. Worse, the experiments were run at
tiny coverage: on NetHEPT ($n = 15233$) a budget of $200$ seeds is $1.3\%$ of the network.

The decisive variable is not the budget but the \emph{seed fraction} $|S|/n$. Saturation is a
property of coverage, not of the number of seeds: the same budget of $200$ seeds is $1.3\%$ of
NetHEPT but $42\%$ of Congress-Twitter ($n = 475$).

\subsection{Sweeping the fraction across eight graphs}

We sweep $|S|/n \in \{0.10, 0.20, 0.40\}$ on eight graphs whose mean degree spans $4.23$ to
$55.95$, and measure the rank correlation between out-degree and the true marginal gain
$\Delta(v \mid S) = \sigma(S \cup \{v\}) - \sigma(S)$. Every cell uses paired Monte Carlo with the
threshold windows shared across candidates within a trial, so the \emph{differences} between
candidates carry no window-draw noise.

Two protocol decisions from \S\ref{sec:protocol} are load-bearing here. The measurement is taken at
$\mathrm{MC} = 300$ paired trials, because at $60$ trials the label noise in the saturated regime is
as large as the dispersion being measured. And each candidate's label carries its own paired
standard error, so the ``neg.'' columns below can separate a candidate that is negative from one
that is merely noisy.

% GENERATED by scripts/audit/write_regime_table.py -- do not edit by hand.
% Source: docs/results/signflip_fixedmodel_mc300.json
% 8 graphs x 3 seed fractions, MC = 300 paired trials per cell.
\begin{table}[t]
\centering
\caption{Out-degree versus the state-conditioned closed form, re-measured under the corrected activation rule at $\mathrm{MC} = 300$ paired trials per cell. Values are means over $|S|/n \ge 20\%$. $\rho$ is the Spearman correlation of the score with the true marginal gain; ``neg.'' is the share of candidates with a negative marginal and ``neg.$^\ast$'' the share negative by more than two paired standard errors. Out-degree is negatively correlated on 7 of 8 graphs and $\delta_2$ is the better ranker on 8 of 8.}
\label{tab:regime-fixed}
\begin{tabular}{lrrrrrr}
\toprule
graph & $n$ & $\langle k \rangle$ & $\rho_{\text{degree}}$ & $\rho_{\delta_2}$ & neg. & neg.$^\ast$ \\
\midrule
NetHEPT            & 15233 & 4.23 & $+0.099$ & $\mathbf{+0.693}$ & 1.0\% & 1.0\% \\
p2p-Gnutella08     & 6301 & 6.59 & $-0.052$ & $\mathbf{+0.400}$ & 0.0\% & 0.0\% \\
ca-GrQc            & 5242 & 11.06 & $-0.583$ & $\mathbf{+0.640}$ & 0.0\% & 0.0\% \\
Wiki-Vote          & 7115 & 29.15 & $-0.042$ & $\mathbf{+0.386}$ & 0.0\% & 0.0\% \\
ca-HepPh           & 12008 & 39.48 & $-0.137$ & $\mathbf{+0.386}$ & 0.0\% & 0.0\% \\
Facebook           & 4039 & 43.69 & $-0.419$ & $\mathbf{+0.527}$ & 2.0\% & 1.0\% \\
email-Eu-core      & 1005 & 50.89 & $-0.236$ & $\mathbf{+0.239}$ & 0.0\% & 0.0\% \\
Congress-Twitter   & 475 & 55.95 & $-0.047$ & $\mathbf{+0.085}$ & 1.0\% & 0.0\% \\
\bottomrule
\end{tabular}
\end{table}


Three observations follow, and the third is a correction of an earlier version of this work.

\begin{enumerate}
  \item \textbf{Out-degree flips sign.} It is positively correlated with the marginal gain on only
        one of the eight graphs --- NetHEPT, the sparsest, at $\langle k \rangle = 4.23$ --- and
        negatively on the other seven once $|S|/n \ge 20\%$. Averaged over the saturated cells the
        correlation is $-0.177$ for out-degree against $+0.420$ for the closed form $\delta_2$.
  \item \textbf{The state-conditioned closed form is the better ranker on all eight.} It is the only
        scorer we tested that is positive in every saturated cell, and its advantage is largest
        exactly where out-degree is most misleading.
  \item \textbf{The objective is almost never non-monotone, and our earlier claim that it often was
        is withdrawn.} An earlier version reported that between $8\%$ and $35\%$ of candidates have
        a negative marginal gain on the denser graphs. Measured at $\mathrm{MC} = 300$ the mean
        share over saturated cells is $0.5\%$, the worst cell in the entire sweep is $4.0\%$, and
        the mean share that is negative by more than two paired standard errors is $0.2\%$. The
        earlier numbers came from twelve trials per cell, where label noise alone produces apparent
        sign changes at that rate; the same measurement at $300$ trials in one of those cells moves
        the share from $44\%$ to $4\%$.
\end{enumerate}

\paragraph{What this means for the argument.} Non-monotonicity is \emph{possible} in this model ---
\S\ref{sec:coverage} exhibits a one-hop instance where adding a seed strictly lowers the target's
activation probability, and that instance is exact arithmetic, not simulation. But it is rare on the
graphs we measure, and the practical difficulty it creates is not that marginal gains go negative.
It is that they become \emph{small} --- near zero and comparable to the error of estimating them ---
while out-degree's association with them reverses. That is why the sign flip is the phenomenon worth
reporting and why $\mathrm{MC} \ge 300$ is a precondition for reporting it at all.

\paragraph{Density is not the explanation.}
It is tempting to attribute the sign flip to density. We do not. Across the eight graphs the rank
correlation between $\langle k \rangle$ and the saturated $\rho_{\text{degree}}$ is $-0.286$, which
is weak, and the relationship is not monotone in density: NetHEPT ($4.23$) is positive,
p2p-Gnutella08 ($6.59$) and Wiki-Vote ($29.15$) are only mildly negative, ca-GrQc ($11.06$) and
Facebook ($43.69$) are strongly negative, and Congress-Twitter ($55.95$), the densest graph, is
nearly zero at $-0.047$. The defensible claim is about the \emph{saturated regime}, not about dense
graphs.

\paragraph{A structural hint.}
After in-weight normalisation, exposure $\delta$ is driven by a node's \emph{in}-degree while a
seed's value is driven by its \emph{out}-degree. The two coincide only when a graph is
approximately symmetric. Consistently, the graphs with $\rho(\text{in}, \text{out})$ near $1$ show
strongly negative out-degree correlations, while NetHEPT, whose in- and out-degrees are nearly
independent, is the one graph where degree survives. We record this as a hypothesis rather than a
result; testing it needs more graphs than the eight available here.
